\section{Conclusions}
% =====================================================================

Nous voulions mesurer quelle information cosmologique les statistiques d'environnement ASTRA ajoutent à la fonction à deux points des galaxies, et comment elles se comparent au density split. Nous avons mesuré les deux sur les mêmes mocks de type DESI et les mêmes halos de Quijote, construit des jeux de covariances et de dérivées pour des prévisions de Fisher, et testé toute la chaîne avec une inférence par émulateur. Les points principaux sont les suivants :
\begin{itemize}
\item \textbf{Gain par rapport à la fonction à deux points.} Les statistiques d'environnement ASTRA contraignent $\omega_b$, $\omega_{\rm cdm}$, $n_s$, $\sigma_8$, $w_0$ et $w_a$ 1.6--3.1 fois mieux que la fonction à deux points des galaxies.
\item \textbf{Comparaison avec le density split.} Les deux méthodes portent une information comparable (rapport d'erreurs 0.8--1.1, le density split légèrement en avance) et sont complémentaires : les combiner gagne encore de 23\% à un facteur 2.4. Avec les halos de Quijote, ASTRA fait mieux sur $\sigma_8$ (facteur 1.5--2) grâce à ses auto-corrélations d'environnement et aux environnements de traceurs.
\item \textbf{Classes de toile cosmique.} Avec les seuils fixes d'ASTRA, les classes se réduisent en pratique à deux classes dominantes (feuillets et filaments) à ces densités, ce qui explique que l'analyse principale utilise des quartiles. Sur Quijote, les classes avec leurs auto-corrélations sont néanmoins la statistique la plus contraignante : deux à trois fois mieux que le density split sur $\sigma_8$, de façon robuste selon les compressions.
\item \textbf{Ce que mesure ASTRA.} Les environnements ASTRA et density split sont fortement mais pas parfaitement corrélés ($\rho\simeq0.72$). La valeur $r$ d'ASTRA elle-même répond à peine à la cosmologie : l'information réside dans la façon dont les environnements se regroupent. ASTRA adapte son échelle à la densité locale, il doit donc être modélisé à la densité et au rapport aléatoires/galaxies des données.
\item \textbf{Méthode.} Les suites d'émulation avec une HOD indépendante par cosmologie ne permettent pas de différences finies naïves, les covariances à boîte unique peuvent biaiser la comparaison des statistiques, et les statistiques d'environnement de traceurs épars ont des dérivées bruitées. Nous avons montré comment chaque problème a été diagnostiqué et traité. L'inférence par l'émulateur est non biaisée et calibrée, mais discriminer les statistiques avec elle exigerait plusieurs réalisations de HOD par point d'entraînement.
\end{itemize}
Les suites naturelles, au-delà du cadre de cet article, sont les classes ASTRA probabilistes de \cite{ASTRA}, qui pourraient rendre utilisables les classes de vides et de nœuds, et un émulateur des statistiques de classes pour une analyse de type DESI des classes.
